%% Caisse glissant à VITESSE CONSTANTE sur un plan incliné à 25° (leçon pc-forces).
%% Principe d'inertie : P + R_N + f = 0. Les trois vecteurs sont tracés à l'échelle
%% (P = 1,6 cm ; R_N = P cos25° = 1,45 cm ; f = P sin25° = 0,68 cm) et se referment,
%% ce que rappelle le petit triangle des forces en haut à gauche.
\documentclass[tikz,border=4pt]{standalone}
\usepackage{liaisons}
\begin{document}
\begin{tikzpicture}[x=1cm,y=1cm,
  vect/.style={-{Stealth[length=6pt]}, line width=1.2pt},
  mince/.style={-{Stealth[length=5pt]}, line width=0.7pt,
                dash pattern=on 3pt off 2pt, black!60}]

\def\alp{25}              % angle d'inclinaison du plan (degrés)
\def\sG{4.41}             % abscisse du centre de la caisse le long de la pente

%% ------------------- le plan incliné (triangle) -----------------------
\filldraw[fill=black!12, draw=black!55, line width=0.8pt, line join=round]
      (0,0) -- (6,0) -- (6,2.798) -- cycle;
% hachures fines sous l'hypoténuse : le « sol »
\begin{scope}[rotate=\alp]
  \foreach \s in {0.9,1.18,...,6.5} { \draw[line width=0.45pt, black!60] (\s,0) -- (\s-0.13,-0.17); }
\end{scope}

%% ------------------- cotation de l'angle ------------------------------
\fill[solideE!25] (0,0) -- (0.9,0) arc (0:\alp:0.9) -- cycle;
\draw[line width=0.5pt, black!55] (0.9,0) arc (0:\alp:0.9);
\node[font=\small] at (12:1.2) {$\alpha$};

%% ------------------- la caisse et son centre G ------------------------
\begin{scope}[rotate=\alp]
  \filldraw[fill=solideD!30, draw=solideD, line width=1pt]
        (\sG-0.45,0) rectangle (\sG+0.45,0.9);
  \coordinate (G) at (\sG,0.45);
  % vitesse constante, vers le bas de la pente, en avant de la caisse
  \draw[mince] (3.72,0.45) -- (2.72,0.45);
  \node[font=\footnotesize, anchor=south east, inner sep=1pt] at (3.68,0.72) {$\vec{v}$ = constante};
\end{scope}
\fill[black] (G) circle (1.6pt);
\node[font=\small, anchor=north east, inner sep=2pt] at (G) {$G$};

%% ------------------- les trois forces ---------------------------------
% poids : vertical, vers le bas
\draw[vect, black] (G) -- ++(0,-1.6)
      node[anchor=west, inner sep=3pt, font=\small] {$\vec{P}$};
% réaction normale : perpendiculaire au plan (25° + 90°)
\draw[vect, solideE] (G) -- ++(\alp+90:1.45)
      node[anchor=south, inner sep=2pt, font=\small, text=solideE, align=center]
      {$\vec{R_N}$\\[-1pt]\footnotesize (réaction normale)};
% frottement : parallèle au plan, vers le HAUT de la pente
\draw[vect, solideB] (G) -- ++(\alp:0.68)
      node[anchor=north west, inner sep=1pt, font=\small, text=solideB] {$\vec{f}$};

%% ------------------- triangle des forces (rappel) ---------------------
\begin{scope}[shift={(0.55,3.40)}]
  \draw[line width=0.7pt, black]   (0,0) -- ++(0,-0.88);
  \draw[line width=0.7pt, solideE] (0,-0.88) -- ++(\alp+90:0.80);
  \draw[line width=0.7pt, solideB] ($(0,-0.88)+(\alp+90:0.80)$) -- ++(\alp:0.374);
\end{scope}
\node[font=\footnotesize, anchor=west] at (0.92,2.98)
     {$\vec{P} + \vec{R_N} + \vec{f} = \vec{0}$};
\end{tikzpicture}
\end{document}
